\documentclass[11pt,a4paper]{article}
\usepackage{operationsform}
\begin{document}
\OPStitle{振り返り・改善提案書}{TPL-DOC-093}{0.2.0}
\OPSsection{1}{基本情報}
\OPSblanklong{案件・文書の識別子}
\OPSblankshort{作成者と連絡先}
\OPSblankshort{作成日と対象版}
\OPSsection{2}{振り返り}
\OPSblanklong{当初の目的と結果}
\OPSblanklong{成功した実務と理由}
\OPSblanklong{困難・不適合の事実}
\OPSblanklong{関係者からの意見}
\OPSsection{3}{改善と継承}
\par\Needspace{4\baselineskip}\noindent\textbf{改善案と優先順位}\par\vspace{1mm}
{\fontsize{8.5}{13}\selectfont
\begin{tabularx}{\linewidth}{@{}XXXX@{}}
\textbf{対象と版} & \textbf{根拠・参照ID} & \textbf{判断・担当} & \textbf{例外と日時}\\[2mm]
\strut & \strut & \strut & \strut\\[5mm]
\strut & \strut & \strut & \strut\\[5mm]
\end{tabularx}
}\vspace{2mm}\par
\OPSblanklong{制度・標準への反映候補}
\OPSblankshort{担当と検証予定}
\OPSblanklong{残された問い}
\OPSsection{4}{確認と履歴}
\par\Needspace{4\baselineskip}\noindent\textbf{参照資料と根拠ID}\par\vspace{1mm}
{\fontsize{8.5}{13}\selectfont
\begin{tabularx}{\linewidth}{@{}XXXX@{}}
\textbf{対象と版} & \textbf{根拠・参照ID} & \textbf{判断・担当} & \textbf{例外と日時}\\[2mm]
\strut & \strut & \strut & \strut\\[5mm]
\strut & \strut & \strut & \strut\\[5mm]
\end{tabularx}
}\vspace{2mm}\par
\OPSblanklong{未確認事項・例外}
\par\Needspace{4\baselineskip}\noindent\textbf{点検・改訂の記録}\par\vspace{1mm}
{\fontsize{8.5}{13}\selectfont
\begin{tabularx}{\linewidth}{@{}XXXX@{}}
\textbf{対象と版} & \textbf{根拠・参照ID} & \textbf{判断・担当} & \textbf{例外と日時}\\[2mm]
\strut & \strut & \strut & \strut\\[5mm]
\strut & \strut & \strut & \strut\\[5mm]
\end{tabularx}
}\vspace{2mm}\par
\end{document}
